% Chapter 8 -- Claim: Simple Trading Rules Beat Buy-and-Hold
% Requires: amsmath, amssymb, booktabs, array, tabularx, tcolorbox (breakable, skins), hyperref or url.
% Defines the callout box only if the preamble has not.
\providecolor{fenavy}{HTML}{1B2A6B}
\providecolor{feshade}{HTML}{F3F5FA}
\makeatletter
% Shaded box for optional technical material that general readers may skip.
\@ifundefined{technote}{%
  \newtcolorbox{technote}{breakable, enhanced,
    colback=feshade, colframe=fenavy!60, boxrule=0.5pt, arc=1mm,
    left=8pt, right=8pt, top=4pt, bottom=4pt}%
}{}
\@ifundefined{callout}{%
  \newtcolorbox{callout}{breakable, enhanced,
    colback=white, colframe=gray!70, boxrule=0.5pt, arc=1mm,
    borderline west={2.5pt}{0pt}{gray!70},
    left=8pt, right=8pt, top=4pt, bottom=4pt}%
}{}
\makeatother

\chapter{Claim: Simple Trading Rules Beat Buy-and-Hold}
\label{ch:08}

``The trend is your friend'' is one of the oldest sayings on Wall Street. It tells a trader to buy what has been rising, sell what has been falling, and stay with the move until it ends. To an efficient-markets economist the saying sounds like superstition: if past prices could predict future prices, traders would use them until the prediction disappeared (Section~\ref{sec:01-efficient-markets}). For decades academic studies agreed, finding that technical trading rules, which are rules based only on past prices and trading volume and are often read from price charts, did not beat buy-and-hold once costs were counted. 

Brock, Lakonishok and LeBaron (1992) reopened the question. They tested two families of simple rules on daily prices of the Dow Jones Industrial Average from 1897 to 1986, about 90 years of data. 
\begin{itemize}
\item \textbf{Moving-average rules}\index{moving-average rule} compare a short moving average of the price, over 1 to 5 days, with a long one, over 50 to 200 days. They buy when the short average rises above the long one and sell when it falls below, so they stay invested while the trend is up and step aside when it turns down (Figure~\ref{fig:06-ma-rule}). 
\item \textbf{Breakout rules}\index{breakout rule} buy when the price rises above its highest level of the past few months, and sell when it falls below its lowest level (Figure~\ref{fig:06-breakout-rule}).
\end{itemize}
Across 26 versions of these rules, the days that followed a buy signal earned much more than the days that followed a sell signal, a gap that standard models of random prices could not reproduce. Many financial websites still draw the 50-day and 200-day moving averages on their stock charts, so the rules they tested are ones that ordinary investors see every day.

\providecolor{fenavy}{HTML}{1B2A6B}
\providecolor{feorange}{HTML}{B8481C}
\begin{figure}[htbp]
\centering
\begin{tikzpicture}[font=\small]
\begin{axis}[width=0.75\linewidth, height=0.36\linewidth, axis lines=left, xmin=0, xmax=600, ymin=70, ymax=210, ylabel={Price}, every tick label/.style={font=\small}, clip=false, title style={font=\small}, legend style={draw=none, fill=none, font=\footnotesize, at={(0.5,-0.32)}, anchor=north, legend columns=-1, /tikz/every even column/.append style={column sep=0.8em}}, legend cell align=left, xlabel={Trading day}]
\fill[fenavy!8] (axis cs:0,70) rectangle (axis cs:221,210);
\fill[fenavy!8] (axis cs:406,70) rectangle (axis cs:600,210);
\addplot[fenavy, thick, no marks] coordinates {(0,125.89) (3,126.85) (6,127.34) (9,127.76) (12,127.60) (15,128.37) (18,129.19) (21,128.79) (24,129.93) (27,132.81) (30,135.63) (33,136.86) (36,137.04) (39,137.31) (42,137.96) (45,137.03) (48,136.85) (51,138.89) (54,140.32) (57,140.66) (60,138.40) (63,137.14) (66,138.82) (69,141.17) (72,144.97) (75,149.47) (78,149.08) (81,148.84) (84,148.29) (87,149.08) (90,149.83) (93,149.56) (96,151.77) (99,154.29) (102,156.25) (105,160.72) (108,163.90) (111,167.36) (114,170.25) (117,171.37) (120,174.22) (123,174.65) (126,174.38) (129,176.58) (132,178.09) (135,180.86) (138,182.89) (141,182.14) (144,185.64) (147,193.07) (150,194.08) (153,191.90) (156,191.52) (159,192.39) (162,190.28) (165,187.74) (168,183.72) (171,180.43) (174,179.54) (177,179.52) (180,178.92) (183,178.64) (186,177.52) (189,175.45) (192,174.58) (195,175.20) (198,177.77) (201,180.65) (204,181.62) (207,181.35) (210,179.37) (213,177.75) (216,175.16) (219,173.80) (222,172.36) (225,168.26) (228,165.19) (231,160.83) (234,156.80) (237,154.65) (240,149.19) (243,144.46) (246,143.17) (249,141.67) (252,138.78) (255,136.63) (258,136.96) (261,135.38) (264,135.76) (267,134.93) (270,134.22) (273,133.76) (276,131.36) (279,129.17) (282,128.95) (285,130.26) (288,130.56) (291,128.50) (294,125.85) (297,122.44) (300,119.52) (303,116.35) (306,112.10) (309,109.02) (312,107.81) (315,107.95) (318,108.73) (321,108.30) (324,107.94) (327,109.64) (330,111.49) (333,110.83) (336,109.62) (339,110.17) (342,109.57) (345,108.18) (348,106.29) (351,105.84) (354,105.03) (357,106.16) (360,108.74) (363,110.09) (366,110.29) (369,113.27) (372,117.53) (375,118.61) (378,118.89) (381,119.60) (384,118.77) (387,118.73) (390,118.99) (393,117.84) (396,116.62) (399,116.57) (402,116.97) (405,117.60) (408,117.72) (411,117.84) (414,117.89) (417,119.32) (420,120.05) (423,121.10) (426,122.76) (429,125.05) (432,126.48) (435,126.07) (438,127.46) (441,131.48) (444,131.74) (447,130.16) (450,130.16) (453,133.13) (456,137.31) (459,138.27) (462,138.64) (465,138.95) (468,138.75) (471,137.52) (474,138.84) (477,141.04) (480,143.29) (483,143.84) (486,142.02) (489,142.24) (492,141.42) (495,139.90) (498,138.31) (501,137.46) (504,138.71) (507,138.71) (510,141.13) (513,143.68) (516,145.34) (519,147.50) (522,145.55) (525,143.21) (528,141.98) (531,140.32) (534,142.07) (537,145.93) (540,147.32) (543,147.89) (546,149.57) (549,151.29) (552,150.27) (555,146.75) (558,146.51) (561,151.34) (564,154.69) (567,156.52) (570,157.84) (573,158.80) (576,156.75) (579,156.27) (582,158.42) (585,160.25) (588,161.89) (591,161.50) (594,161.59) (597,164.29) (600,165.20)};
\addlegendentry{Short average (5 days)}
\addplot[black, very thick, dashed, no marks] coordinates {(0,107.68) (3,108.23) (6,108.75) (9,109.32) (12,109.86) (15,110.42) (18,110.94) (21,111.50) (24,112.08) (27,112.72) (30,113.44) (33,114.21) (36,114.95) (39,115.70) (42,116.44) (45,117.16) (48,117.91) (51,118.67) (54,119.46) (57,120.25) (60,120.97) (63,121.63) (66,122.31) (69,123.04) (72,123.85) (75,124.71) (78,125.51) (81,126.40) (84,127.25) (87,128.12) (90,129.00) (93,129.94) (96,130.94) (99,131.95) (102,132.97) (105,134.06) (108,135.16) (111,136.35) (114,137.54) (117,138.72) (120,140.00) (123,141.21) (126,142.35) (129,143.48) (132,144.56) (135,145.67) (138,146.84) (141,148.01) (144,149.26) (147,150.63) (150,151.99) (153,153.27) (156,154.57) (159,155.85) (162,157.09) (165,158.24) (168,159.31) (171,160.32) (174,161.31) (177,162.22) (180,163.07) (183,163.89) (186,164.69) (189,165.43) (192,166.15) (195,166.93) (198,167.76) (201,168.61) (204,169.42) (207,170.25) (210,171.05) (213,171.86) (216,172.53) (219,173.18) (222,173.67) (225,174.02) (228,174.33) (231,174.52) (234,174.69) (237,174.77) (240,174.71) (243,174.60) (246,174.39) (249,174.11) (252,173.71) (255,173.21) (258,172.65) (261,171.97) (264,171.28) (267,170.53) (270,169.70) (273,168.89) (276,168.01) (279,167.02) (282,166.04) (285,165.01) (288,163.94) (291,162.87) (294,161.59) (297,160.12) (300,158.62) (303,157.10) (306,155.47) (309,153.79) (312,152.16) (315,150.58) (318,149.11) (321,147.67) (324,146.23) (327,144.85) (330,143.52) (333,142.14) (336,140.81) (339,139.51) (342,138.21) (345,136.84) (348,135.39) (351,133.87) (354,132.34) (357,130.85) (360,129.48) (363,128.14) (366,126.87) (369,125.68) (372,124.63) (375,123.64) (378,122.75) (381,121.96) (384,121.20) (387,120.51) (390,119.95) (393,119.41) (396,118.89) (399,118.39) (402,118.00) (405,117.61) (408,117.24) (411,116.90) (414,116.53) (417,116.24) (420,115.97) (423,115.73) (426,115.59) (429,115.53) (432,115.49) (435,115.37) (438,115.35) (441,115.45) (444,115.57) (447,115.74) (450,115.98) (453,116.37) (456,116.92) (459,117.51) (462,118.13) (465,118.76) (468,119.35) (471,119.93) (474,120.56) (477,121.19) (480,121.84) (483,122.51) (486,123.14) (489,123.79) (492,124.42) (495,125.05) (498,125.69) (501,126.32) (504,127.01) (507,127.64) (510,128.30) (513,128.97) (516,129.68) (519,130.34) (522,130.86) (525,131.35) (528,131.79) (531,132.20) (534,132.72) (537,133.27) (540,133.85) (543,134.47) (546,135.14) (549,135.83) (552,136.47) (555,137.02) (558,137.62) (561,138.33) (564,139.07) (567,139.84) (570,140.58) (573,141.33) (576,141.98) (579,142.60) (582,143.25) (585,143.96) (588,144.62) (591,145.19) (594,145.82) (597,146.53) (600,147.21)};
\addlegendentry{Long average (150 days)}
\addlegendimage{area legend, fill=fenavy!8, draw=fenavy!8}
\addlegendentry{Rule holds the asset}
\fill[feorange] (axis cs:222,173.67) circle (2.2pt);
\draw[feorange, thin] (axis cs:222,173.67) -- (axis cs:247,189.7) node[anchor=south west, inner sep=1pt, font=\footnotesize] {sell: short crosses below long};
\fill[fenavy] (axis cs:406,117.48) circle (2.2pt);
\draw[fenavy, thin] (axis cs:406,117.48) -- (axis cs:431,103.5) node[anchor=north west, inner sep=1pt, font=\footnotesize] {buy: short crosses above long};
\end{axis}
\end{tikzpicture}
\caption{A moving-average rule stays invested while the trend is up and steps aside when it turns down, but always a little late. The moving-average rule applied to a simulated price. The rule holds the asset (shaded) while a short average of the price (5 days) is above a long one (150 days). It rides the long rise and avoids most of the fall, but it acts only after each turn has begun.}
\label{fig:06-ma-rule}
\end{figure}

\begin{figure}[htbp]
\centering
\begin{tikzpicture}[font=\small]
\begin{axis}[width=0.75\linewidth, height=0.36\linewidth, axis lines=left, xmin=0, xmax=600, ymin=70, ymax=210, ylabel={Price}, every tick label/.style={font=\small}, clip=false, title style={font=\small}, legend style={draw=none, fill=none, font=\footnotesize, at={(0.5,-0.32)}, anchor=north, legend columns=-1, /tikz/every even column/.append style={column sep=0.8em}}, legend cell align=left, xlabel={Trading day}]
\fill[fenavy!8] (axis cs:0,70) rectangle (axis cs:180,210);
\fill[fenavy!8] (axis cs:368,70) rectangle (axis cs:600,210);
\addplot[fenavy, thick, no marks] coordinates {(0,126.42) (3,127.99) (6,127.71) (9,128.33) (12,126.99) (15,131.02) (18,127.35) (21,129.39) (24,131.76) (27,135.64) (30,136.62) (33,138.24) (36,136.88) (39,138.13) (42,137.59) (45,136.02) (48,138.94) (51,138.85) (54,140.35) (57,139.00) (60,136.06) (63,137.35) (66,141.08) (69,142.97) (72,149.90) (75,149.25) (78,148.52) (81,149.74) (84,148.41) (87,151.34) (90,147.76) (93,151.15) (96,153.18) (99,155.13) (102,158.58) (105,164.27) (108,165.14) (111,169.10) (114,170.97) (117,174.46) (120,175.24) (123,173.70) (126,175.18) (129,178.76) (132,179.23) (135,182.35) (138,183.04) (141,182.35) (144,192.02) (147,194.59) (150,194.64) (153,189.77) (156,192.41) (159,190.51) (162,187.25) (165,186.07) (168,180.63) (171,178.56) (174,179.30) (177,179.54) (180,178.95) (183,178.86) (186,175.13) (189,176.43) (192,174.84) (195,177.91) (198,178.49) (201,181.46) (204,181.54) (207,180.92) (210,178.04) (213,176.77) (216,171.36) (219,173.75) (222,169.02) (225,167.34) (228,164.58) (231,157.44) (234,155.83) (237,152.70) (240,142.85) (243,145.59) (246,142.03) (249,141.34) (252,134.61) (255,138.59) (258,135.06) (261,135.09) (264,136.18) (267,134.02) (270,135.19) (273,131.77) (276,129.26) (279,128.08) (282,128.22) (285,130.91) (288,129.27) (291,126.99) (294,123.45) (297,120.03) (300,116.62) (303,113.64) (306,109.31) (309,108.39) (312,108.37) (315,109.57) (318,108.68) (321,107.22) (324,108.24) (327,111.19) (330,113.36) (333,108.61) (336,110.10) (339,109.18) (342,109.15) (345,106.23) (348,105.62) (351,104.92) (354,104.77) (357,108.22) (360,110.41) (363,109.87) (366,110.61) (369,117.14) (372,118.98) (375,118.52) (378,119.44) (381,119.50) (384,118.33) (387,119.23) (390,118.15) (393,116.51) (396,115.08) (399,116.63) (402,117.24) (405,117.48) (408,116.76) (411,116.48) (414,120.12) (417,118.88) (420,120.96) (423,121.56) (426,125.37) (429,124.48) (432,127.79) (435,125.10) (438,130.36) (441,132.53) (444,130.31) (447,129.45) (450,132.02) (453,137.77) (456,139.12) (459,138.54) (462,138.59) (465,138.89) (468,138.61) (471,138.06) (474,138.99) (477,142.11) (480,145.30) (483,141.03) (486,143.13) (489,141.81) (492,141.46) (495,138.04) (498,137.08) (501,138.23) (504,138.73) (507,140.21) (510,143.03) (513,144.00) (516,146.40) (519,147.78) (522,143.18) (525,143.49) (528,140.65) (531,139.21) (534,145.47) (537,146.66) (540,147.31) (543,149.25) (546,150.80) (549,151.07) (552,146.71) (555,144.92) (558,149.59) (561,154.09) (564,154.80) (567,156.91) (570,159.22) (573,156.93) (576,155.79) (579,157.00) (582,160.28) (585,161.65) (588,160.92) (591,159.74) (594,162.72) (597,164.80) (600,164.51)};
\addlegendentry{Daily price}
\addplot[black, thick, dashed, no marks] coordinates {(0,127.93) (3,127.93) (6,127.99) (9,127.99) (12,128.33) (15,128.39) (18,131.45) (21,131.45) (24,131.45) (27,133.54) (30,136.92) (33,137.03) (36,138.45) (39,138.48) (42,138.61) (45,138.61) (48,138.61) (51,139.91) (54,142.47) (57,142.47) (60,142.47) (63,142.47) (66,142.47) (69,142.47) (72,146.45) (75,150.90) (78,150.90) (81,150.90) (84,150.90) (87,150.90) (90,151.34) (93,151.34) (96,152.75) (99,155.22) (102,156.91) (105,161.97) (108,164.68) (111,169.59) (114,171.23) (117,171.23) (120,177.00) (123,177.00) (126,177.00) (129,178.16) (132,178.76) (135,182.68) (138,184.58) (141,184.58) (144,189.07) (147,194.85) (150,195.19) (153,195.19) (156,195.19) (159,195.19) (162,195.19) (165,195.19) (168,195.19) (171,195.19) (174,195.19) (177,195.19) (180,195.19) (183,195.19) (186,195.19) (189,195.19) (192,195.19) (195,195.19) (198,195.19) (201,193.37) (204,193.37) (207,193.01) (210,190.59) (213,189.08) (216,183.88) (219,183.05) (222,183.05) (225,183.05) (228,183.05) (231,183.05) (234,183.05) (237,183.05) (240,183.05) (243,183.05) (246,183.05) (249,183.05) (252,182.49) (255,182.49) (258,178.52) (261,178.22) (264,175.46) (267,175.26) (270,175.04) (273,168.76) (276,164.58) (279,160.66) (282,157.05) (285,155.79) (288,150.32) (291,145.59) (294,143.00) (297,141.60) (300,140.01) (303,138.59) (306,137.35) (309,137.09) (312,137.09) (315,135.19) (318,135.19) (321,134.76) (324,132.06) (327,132.06) (330,132.06) (333,132.06) (336,131.35) (339,128.53) (342,125.60) (345,122.33) (348,119.71) (351,116.50) (354,113.36) (357,113.36) (360,113.36) (363,113.36) (366,113.36) (369,115.74) (372,118.49) (375,118.98) (378,119.87) (381,120.12) (384,120.12) (387,120.12) (390,120.12) (393,120.12) (396,120.12) (399,120.12) (402,120.12) (405,120.12) (408,120.12) (411,120.12) (414,120.12) (417,120.69) (420,121.11) (423,121.57) (426,123.07) (429,126.46) (432,127.54) (435,127.79) (438,129.05) (441,133.76) (444,133.76) (447,133.76) (450,133.76) (453,135.26) (456,138.00) (459,139.12) (462,139.30) (465,139.46) (468,139.46) (471,139.46) (474,140.46) (477,142.48) (480,145.03) (483,145.30) (486,145.30) (489,145.30) (492,145.30) (495,145.30) (498,145.30) (501,145.30) (504,145.30) (507,145.30) (510,145.30) (513,145.30) (516,145.93) (519,149.22) (522,149.22) (525,149.22) (528,149.22) (531,149.22) (534,149.22) (537,149.22) (540,149.22) (543,149.22) (546,152.08) (549,152.09) (552,152.09) (555,152.09) (558,152.09) (561,153.88) (564,155.98) (567,158.86) (570,158.86) (573,161.24) (576,161.24) (579,161.24) (582,161.24) (585,161.24) (588,163.24) (591,163.24) (594,163.24) (597,167.52) (600,167.52)};
\addlegendentry{50-day high}
\addplot[black, thick, dotted, no marks] coordinates {(0,105.83) (3,107.14) (6,108.46) (9,108.46) (12,110.62) (15,111.09) (18,111.09) (21,113.02) (24,115.93) (27,119.77) (30,122.41) (33,122.41) (36,122.41) (39,122.41) (42,124.66) (45,125.40) (48,125.40) (51,126.78) (54,126.99) (57,126.99) (60,126.99) (63,127.35) (66,127.35) (69,128.42) (72,128.42) (75,132.81) (78,134.61) (81,134.61) (84,134.61) (87,135.26) (90,135.26) (93,135.26) (96,135.26) (99,136.06) (102,136.06) (105,136.06) (108,136.06) (111,136.47) (114,138.46) (117,141.42) (120,143.60) (123,146.58) (126,146.58) (129,146.58) (132,146.58) (135,147.76) (138,147.76) (141,149.20) (144,152.41) (147,155.13) (150,155.44) (153,161.87) (156,163.44) (159,168.29) (162,169.43) (165,169.43) (168,173.64) (171,173.64) (174,173.95) (177,176.84) (180,177.36) (183,177.03) (186,176.21) (189,174.71) (192,173.01) (195,173.01) (198,173.01) (201,173.01) (204,173.01) (207,173.01) (210,173.01) (213,173.01) (216,173.01) (219,171.36) (222,169.34) (225,166.85) (228,163.43) (231,158.03) (234,155.66) (237,153.28) (240,146.79) (243,142.85) (246,142.02) (249,141.36) (252,136.34) (255,134.61) (258,134.61) (261,134.61) (264,134.61) (267,133.90) (270,133.35) (273,133.35) (276,130.28) (279,128.86) (282,128.08) (285,128.08) (288,128.08) (291,128.08) (294,124.68) (297,121.73) (300,119.49) (303,115.28) (306,110.48) (309,107.61) (312,107.16) (315,106.72) (318,106.72) (321,106.72) (324,106.72) (327,106.72) (330,106.72) (333,106.72) (336,106.72) (339,106.72) (342,106.72) (345,106.72) (348,106.07) (351,105.62) (354,104.73) (357,104.73) (360,104.73) (363,104.73) (366,104.73) (369,104.73) (372,104.73) (375,104.73) (378,104.73) (381,104.73) (384,104.73) (387,104.73) (390,104.73) (393,104.73) (396,104.73) (399,104.73) (402,104.73) (405,106.18) (408,109.08) (411,109.87) (414,110.00) (417,112.83) (420,115.08) (423,115.08) (426,115.08) (429,115.08) (432,115.08) (435,115.08) (438,115.08) (441,115.08) (444,115.08) (447,115.75) (450,115.75) (453,115.75) (456,115.75) (459,115.75) (462,115.75) (465,118.88) (468,120.32) (471,120.32) (474,122.21) (477,124.48) (480,124.74) (483,124.74) (486,128.02) (489,128.99) (492,128.99) (495,128.99) (498,128.99) (501,131.59) (504,135.37) (507,135.37) (510,135.37) (513,135.37) (516,135.37) (519,135.37) (522,136.54) (525,136.54) (528,136.54) (531,136.54) (534,136.54) (537,136.54) (540,136.54) (543,136.54) (546,136.54) (549,136.54) (552,136.72) (555,136.72) (558,139.21) (561,139.21) (564,139.21) (567,139.21) (570,139.21) (573,139.21) (576,139.21) (579,139.21) (582,141.93) (585,144.92) (588,144.92) (591,144.92) (594,144.92) (597,144.92) (600,144.92)};
\addlegendentry{50-day low}
\fill[feorange] (axis cs:181,177.03) circle (2.2pt);
\draw[feorange, thin] (axis cs:181,177.03) -- (axis cs:206,199.0) node[anchor=south west, inner sep=1pt, font=\footnotesize] {sell: breaks below 50-day low};
\fill[fenavy] (axis cs:368,115.74) circle (2.2pt);
\draw[fenavy, thin] (axis cs:368,115.74) -- (axis cs:393,93.7) node[anchor=north west, inner sep=1pt, font=\footnotesize] {buy: breaks above 50-day high};
\end{axis}
\end{tikzpicture}
\caption{A breakout rule buys new highs and sells new lows. The breakout rule applied to the same simulated price as in Figure~\ref{fig:06-ma-rule}. The rule buys when the price rises above its highest level of the previous 50 days and sells when it falls below its lowest level; shaded periods are the days it holds the asset. Like the moving-average rule, it acts only after each turn has begun, but in this example it reacts sooner.}
\label{fig:06-breakout-rule}
\end{figure}


The finding mattered because it challenged the weak form of market efficiency (Section~\ref{sec:01-efficient-markets}), the idea that past prices cannot predict future returns, with rules anyone could follow, and because a large industry already traded on such rules. It also raised a methodological problem that became central to empirical finance. Trend rules come in thousands of variants, and any of them can be tuned on the past. If the rules that Brock, Lakonishok and LeBaron tested had survived because generations of traders kept the ones that worked on the Dow, their success might be nothing more than the best result of a long search. Sullivan, Timmermann and White (1999) set out to measure that risk. They applied White's Reality Check, a test that accounts for the whole universe of rules a researcher might have tried, and found that the best rule's success on the Dow did not carry into the following decade. Moskowitz, Ooi and Pedersen (2012) then asked whether trends were a quirk of one stock index or a feature of markets in general, and they documented \emph{time-series momentum} in 58 futures markets. Their strongest evidence came from a test that pooled all the markets together, however, and Huang, Li, Wang and Zhou (2020) questioned how strong that evidence is once each market is tested on its own.

The implications are practical. Managed-futures funds, often called commodity trading advisors, which are professional funds that trade futures contracts in many markets, run large sums on trend rules and market them as protection in crises, so whether trends are real decides whether their investors are paid for anything. For research, the chapter shows how to judge a claim that was found by searching. The claim on trial\index{claims-on-trial method} is: \textbf{simple trend rules on an asset's own past prices earn returns that survive a full accounting for the number of rules tried.}

The lab in Section~\ref{sec:08-lab} replicates the 26 rules of Brock, Lakonishok and LeBaron (1992) on the free daily U.S. market series of the Kenneth R. French Data Library since 1926, applies the Reality Check of Sullivan, Timmermann and White (1999) to the full rule universe, and rebuilds the time-series momentum of Moskowitz, Ooi and Pedersen (2012) with AQR's free factor data and FRED series, in order to check whether trend signals still predict returns after data snooping, costs and publication.

By the end of the chapter you will be able to:

\begin{itemize}
  \item define the moving-average and breakout rules and compute their signals;
  \item contrast time-series momentum\index{momentum!time-series} with the cross-sectional momentum of Chapter~\ref{ch:05};
  \item size positions to a volatility target\index{volatility targeting} using only past data;
  \item test the best of many rules with White's Reality Check\index{Reality Check} and Hansen's SPA test\index{SPA test};
  \item explain why trend following resembles owning a straddle\index{straddle};
  \item test the claim after publication, after costs and asset by asset.
\end{itemize}

\begin{callout}
\textbf{How to read this chapter.} The chapter is written for two kinds of readers at once. If you are new to finance, you can follow the whole argument by reading the main text, the examples and the short \emph{In plain terms} summaries; the box ``For readers new to the stock market'' in Chapter~\ref{ch:01} explains the basic market terms used here. The shaded \emph{Technical note} boxes contain statistical details, and you can skip them without losing the thread. The Lab, the robustness checks and the lab extensions are written for readers who want to run the Python code themselves.
\end{callout}

\section{Foundations: technical rules and time-series momentum}
\label{sec:08-foundations}

\subsection{Moving-average and trading-range breakout rules}

A \textbf{moving average}\index{moving average|textbf} is the average of the most recent prices. Write $P_t$ for the price (or index level) on day $t$. The $n$-day moving average is
\begin{equation}
  MA_t(n) = \frac{1}{n} \sum_{i=0}^{n-1} P_{t-i}.
  \label{eq:08-ma}
\end{equation}
A one-day moving average is just today's price. A long moving average, such as 200 days, moves slowly and shows the general direction of the market. A short one, such as 1 or 5 days, follows the price closely.

A \textbf{moving-average rule}\index{technical trading rule!moving average|textbf} compares a short average with a long one. It gives a buy signal when the short average is above the long one by more than a band $b$, and a sell signal when it is below by more than the band:
\begin{equation}
  \text{buy if } MA_t(s) > (1+b)\,MA_t(l), \qquad \text{sell if } MA_t(s) < (1-b)\,MA_t(l),
  \label{eq:08-marule}
\end{equation}
where $s$ is the short window and $l$ the long one. Inside the band no new signal is given. The band filters out ``whipsaws,'' the many small crossings that happen when the two averages are close. A rule is described by its two windows and its band, so ``1--200, 1\%'' means today's price against the 200-day average with a 1\% band.

Brock, Lakonishok and LeBaron used the rule in two ways. In the \textbf{variable-length moving average} (VMA) version, the investor holds the position for as long as the signal lasts. In the \textbf{fixed-length moving average} (FMA) version, the investor acts only when the short average crosses the long one and then holds for a fixed ten days, which lets the researcher measure what happens after a crossing.

A \textbf{trading-range breakout rule}\index{technical trading rule!trading-range breakout|textbf} uses recent highs and lows instead of averages. Traders call the recent high a \textbf{resistance level} and the recent low a \textbf{support level}. The rule buys when the price rises above the maximum of the past $n$ days and sells when it falls below the minimum:
\begin{equation}
  \text{buy if } P_t > (1+b)\max_{1\le i\le n} P_{t-i}, \qquad \text{sell if } P_t < (1-b)\min_{1\le i\le n} P_{t-i}.
  \label{eq:08-trb}
\end{equation}

\textbf{Worked example.} Suppose an index closes at 104 today. Its 200-day average is 100 and its 50-day high is 103. The 1--200 rule with a 1\% band compares 104 with $1.01 \times 100 = 101$ and signals a buy. The 50-day breakout rule with a 1\% band compares 104 with $1.01 \times 103 = 104.03$ and gives no signal, because the price has not cleared the band.

\subsection{Time-series versus cross-sectional momentum}

Chapter~\ref{ch:05} studied \textbf{cross-sectional momentum}: it ranks stocks against each other and buys the relative winners. \textbf{Time-series momentum}\index{momentum!time-series|textbf} looks at each asset on its own. It asks only whether the asset's own excess return over the look-back period was positive or negative, and holds it long or short accordingly. In a year when every market rises, cross-sectional momentum is still long half the assets and short the other half, while time-series momentum is long all of them.

Moskowitz, Ooi and Pedersen define the strategy for asset $s$ as follows. At the end of month $t$, look at the asset's excess return over the past twelve months, $r^s_{t-12,t}$. If it is positive, go long for one month; if it is negative, go short. The position is scaled by the asset's volatility, as described in the next subsection. The monthly return is
\begin{equation}
  r^{\text{TSMOM},s}_{t,t+1} = \operatorname{sign}\left(r^s_{t-12,t}\right) \frac{40\%}{\sigma^s_t}\, r^s_{t,t+1},
  \label{eq:08-tsmom}
\end{equation}
where $\sigma^s_t$ is the asset's annualized volatility estimated with data up to month $t$. A twelve-month time-series momentum signal is a close cousin of a moving-average rule: an asset whose price is above its level of a year ago is usually also above its long moving average. Both rules ask the same question, ``has the price been going up?'', in slightly different ways.

\subsection{Volatility targeting across assets}

A trend-following portfolio may hold a Treasury bond, whose volatility is about 5\% a year, next to crude oil, whose volatility can exceed 40\%. If both received the same dollar amount, oil would dominate the portfolio's risk. \textbf{Volatility targeting}\index{volatility targeting|textbf} sizes each position so that each contributes about the same risk. In plain words, a calm market such as government bonds gets a large position and a wild one such as oil gets a small one, so that neither dominates. The shaded box shows how the risk of each market is estimated.

\begin{technote}
\textbf{Technical note: estimating each market's volatility.} Moskowitz, Ooi and Pedersen estimate each asset's variance with an exponentially weighted average of past squared daily returns,
\begin{equation}
  \left(\sigma^s_t\right)^2 = 261 \sum_{i=0}^{\infty} (1-\delta)\,\delta^i \left(r^s_{t-1-i} - \bar r^s_t\right)^2,
  \label{eq:08-ewma}
\end{equation}
where 261 annualizes a daily variance and $\delta$ is chosen so that the average age of the data, $\delta/(1-\delta)$, is 60 days. This is the EWMA estimator that Chapter~\ref{ch:10} develops further. The position in Equation~\eqref{eq:08-tsmom} is then $40\%/\sigma^s_t$, which targets an annual volatility of 40\% for each asset. With dozens of assets that do not move together, the diversified portfolio's volatility is far lower, about 12\% a year in their sample.
\end{technote}

Two rules keep the test honest. The volatility estimate must use only data available before the position is taken, and the target must be fixed in advance, not chosen after seeing which value worked best.

\subsection{Rule universes and the SPA test}

A single rule tested once is one hypothesis. Technical analysis offers thousands: every pair of window lengths, every band, every holding period and every family of rules. Sullivan, Timmermann and White (1999) built a universe of 7,846 rules from five families: filter rules, moving averages, support and resistance, channel breakouts and on-balance volume. The best rule in such a universe will look good by chance alone, exactly as the best of 1,000 random strategies did in Lab~\ref{ch:01}.

White's Reality Check (Section~\ref{sec:09-reality}) answers the right question: could the best rule have done this well if no rule had any real ability? The answer is reported as a $p$-value, the probability that luck alone would produce a best rule this good, so a small $p$-value is evidence of real ability. The shaded box explains how the test and its successor work.

\begin{technote}
\textbf{Technical note: the Reality Check and the SPA test.} The Reality Check compares the best rule's performance with the distribution of the \emph{best} performance across all rules in bootstrap samples where the true performance of every rule is set to zero. The bootstrap must keep the dependence in daily returns, so it resamples blocks of days; Sullivan, Timmermann and White used the stationary bootstrap of Politis and Romano (1994), with blocks of random length averaging ten days.

The Reality Check has one weakness. Adding many hopeless rules to the universe makes it harder to declare the best rule significant, even though the hopeless rules are irrelevant. Hansen's (2005) \textbf{Superior Predictive Ability} (SPA) test\index{SPA test|textbf} fixes this. It divides each rule's average performance by its own standard error, and it largely ignores rules that clearly do worse than the benchmark when it builds the null distribution. The SPA test therefore has more power when a universe contains many poor rules, which is the usual case.
\end{technote}

\noindent\textbf{In plain terms.} Trend rules are simple: own an asset while its price has been rising and get out, or bet against it, when the price has been falling. Because there are thousands of ways to write such a rule, the best-looking rule in a backtest is bound to look good partly by luck, and special tests are needed to judge whether it is better than that.

\section{How it works}
\label{sec:08-how}

\textbf{Why trends might exist.} Several explanations have been proposed, and they are not mutually exclusive.

\begin{itemize}
  \item \textbf{Slow reaction to news.} Investors anchor on old prices and update their views gradually, so prices react to news too little at first and keep moving as the news sinks in. This is the underreaction\index{underreaction} story of Chapter~\ref{ch:05}.
  \item \textbf{Herding and feedback.} Once a move is under way, other investors chase it, and the move overshoots. This explains why trends eventually partly reverse, as Moskowitz, Ooi and Pedersen found after about a year.
  \item \textbf{Investors who do not trade for profit.} Central banks smooth exchange rates and interest rates, and hedgers such as airlines and farmers trade to reduce risk rather than to earn a return. Their trades can slow the adjustment of prices. Moskowitz, Ooi and Pedersen found that speculators in futures markets tend to follow trends, and that they profit at the expense of hedgers.
\end{itemize}

\textbf{Why trend following likes extreme markets.} A trend rule is long after prices have risen and short after they have fallen. In a long, steady decline, such as 2008 or the bear markets of the 1970s, it turns short early and profits as the decline continues. In a long, steady rise it is long. It loses in markets that move back and forth, when every signal arrives just before the price turns. The payoff therefore looks like that of an option \textbf{straddle}\index{straddle|textbf}, a combination of a call and a put that profits from a large move in either direction. A call option pays off when a price rises past a set level, and a put pays off when it falls past one, so owning both pays when prices move a lot in either direction. Fung and Hsieh (2001) first described trend followers' returns in these terms. Moskowitz, Ooi and Pedersen confirmed the pattern: regressing the strategy's quarterly return on the market's return and its square gave an insignificant linear term and a significant positive square term, a ``smile.''

The straddle pattern explains the appeal of trend following as a \textbf{diversifier}\index{diversification}. It tends to do well when a stock portfolio does badly over several months. The pattern has a limit: a straddle profits from large moves that last long enough for the rule to react. A sudden one-day crash, such as October 1987, or a sharp turn, such as March 2009, can hurt a trend follower who was positioned the other way.

\noindent\textbf{In plain terms.} Trend following tends to profit from long, steady moves in either direction, including long market declines, and to lose money when markets chop back and forth. That pattern is why many investors hold it as a form of insurance for a stock portfolio, while knowing it cannot protect against a sudden one-day crash.

\section{The claim}
\label{sec:08-claim}

Stated so that it can fail:

\begin{callout}
\textbf{Simple trend rules, which buy assets whose prices have been rising and sell assets whose prices have been falling, earn returns above a buy-and-hold benchmark, and the best of them remain significant after accounting for the number of rules tried, after publication, and after trading costs.}
\end{callout}

The claim has two forms. The \textbf{strong form} is that classic technical rules on a single stock index, such as the Dow, beat buy-and-hold. The \textbf{weak form} is that a diversified time-series momentum portfolio across many asset classes earns a positive risk-adjusted return. The tests follow:

\begin{itemize}
  \item \textbf{Search.} Is the best rule still significant after the Reality Check and the SPA test?
  \item \textbf{Time.} Did the rules keep working after the samples that established them, and after they were published?
  \item \textbf{Breadth.} Is the evidence strong asset by asset, or only in a pooled test?
  \item \textbf{Costs.} Trend rules trade often. Does the profit survive realistic costs?
\end{itemize}

\section{The original evidence}
\label{sec:08-evidence}

\subsection{Brock, Lakonishok and LeBaron (1992)}

Brock, Lakonishok and LeBaron\index{Brock, Lakonishok and LeBaron|textbf} tested 26 rules on the daily Dow Jones Industrial Average from 1897 to 1986. Ten were variable-length moving-average rules (the window pairs 1--50, 1--150, 5--150, 1--200 and 2--200, each with no band and with a 1\% band), ten were the same rules in fixed-length form with a ten-day holding period, and six were trading-range breakout rules with windows of 50, 150 and 200 days, with and without a 1\% band.

The results were striking for their time. Returns after buy signals were consistently higher than returns after sell signals. For the fixed-length rules, the average ten-day return after a buy signal was about 0.53\%, the average after a sell signal was about $-0.40\%$, and the average ten-day return over the whole sample was 0.17\%. Returns after sell signals were also more volatile. A negative expected return after sell signals is hard to explain with any model in which investors are paid for risk.

The authors used a bootstrap to check the results. They fitted four standard models of returns to the data, a random walk, an AR(1), a GARCH-in-mean and an exponential GARCH, which range from purely random price changes to models in which volatility rises and falls over time, simulated many artificial Dow histories from each, and applied the rules to the simulated data. None of the four models produced the buy--sell differences seen in the real data.

\subsection{Sullivan, Timmermann and White (1999)}

Sullivan, Timmermann and White asked the question that Brock, Lakonishok and LeBaron could not: how many rules had the profession tried before settling on these 26? They extended the Dow data to 1996, built the universe of 7,846 rules described in Section~\ref{sec:08-foundations}, and applied White's Reality Check.

In the original 1897--1986 sample the answer favored technical analysis. The best of Brock, Lakonishok and LeBaron's own rules, a 50-day variable moving average with a 1\% band, earned about 10\% a year, and the best rule in the full universe, a 5-day moving average, earned about 19\% a year. Both had Reality Check $p$-values close to zero.

The ten years after the original sample told a different story. From 1987 to 1996 the best rule in the full universe had a $p$-value of 0.34, and the best of the original 26 had a $p$-value of 0.15. On S\&P 500 futures from 1984 to 1996, no rule beat the benchmark. The rules had worked in the sample in which they were discovered and not in the decade that followed.

\subsection{Moskowitz, Ooi and Pedersen (2012) and Hurst, Ooi and Pedersen (2017)}

Moskowitz, Ooi and Pedersen\index{Moskowitz, Ooi and Pedersen|textbf} moved the question from one stock index to many markets. They studied 58 liquid futures and forward contracts, meaning contracts that trade in large volumes and can be bought and sold cheaply: 24 commodities, 12 currency pairs, 9 equity indices and 13 government bond futures. Their price data begin in 1965, and they evaluated the strategy from 1985 to 2009, when many more markets were available.

Three findings stand out. First, in a pooled regression of each asset's volatility-scaled monthly return on its own past returns, the coefficients were positive for the first twelve months of past returns and mostly negative beyond, so trends persisted for about a year and then partly reversed. Second, the twelve-month strategy of Equation~\eqref{eq:08-tsmom} had a positive average return in every one of the 58 contracts, and 52 were significant at the 5\% level. Third, the diversified portfolio had a Sharpe ratio above one and an alpha of about 1.6\% a month against a global stock index and the size, value and momentum factors. It performed best in extreme market quarters, up or down.

Hurst, Ooi and Pedersen (2017) pushed the test back in time. They combined one-, three- and twelve-month signals in 67 markets, including commodities, equity indices, bond markets and currency pairs, and simulated the strategy from 1880 to 2016, with higher assumed trading costs in the early decades and after a 2\% management fee and a 20\% performance fee, the typical fees of a hedge fund, which take 2\% of the money each year plus a fifth of any profit. The strategy earned positive returns in every decade, had low correlations with stocks and bonds, and made money in most of the ten largest drawdowns of a 60/40 stock--bond portfolio, the classic mix of 60\% stocks and 40\% bonds, where a drawdown is a fall from a previous peak.

\subsection{Huang, Li, Wang and Zhou (2020)}

Huang, Li, Wang and Zhou\index{time-series momentum!critique} challenged the statistical evidence. Their argument has three parts.

\begin{itemize}
  \item \textbf{The pooled regression overstates the evidence.} A regression that stacks all assets together can find a significant slope even when no single asset is predictable, for example when assets with higher average returns also tend to have positive past returns. When Huang and coauthors bootstrapped the pooled $t$-statistic, it was no longer significant.
  \item \textbf{Asset by asset, the evidence is weak.} Running the predictive regression separately for each of their 55 assets, most slopes were insignificant, and only a few assets showed significant out-of-sample predictability.
  \item \textbf{A simpler strategy does as well.} A ``time-series history'' strategy, which goes long an asset whose average return over its whole past is positive and short otherwise, earned about the same as time-series momentum. That strategy needs no predictability at all. It profits from the fact that assets with positive average returns in the past, such as stock indices and bonds, kept earning positive average returns.
\end{itemize}

Their conclusion was that the strategy is profitable but that the profit need not come from trends. Defenders of trend following reply that the strategy's value lies in its behavior in crises, not only in its average return, and that a diversified portfolio of weak signals can be valuable even when each signal is hard to detect on its own (Section~\ref{sec:08-stress}).

\noindent\textbf{In plain terms.} Simple trend rules on the Dow looked impressive for most of the twentieth century but stopped working once they were tested on the years after their discovery. Trend following across dozens of markets has a better record and has often made money in long market declines, but critics argue that part of its profit comes simply from holding assets that tend to rise over time.

\section{Lab~\thechapter: Replicating Brock--Lakonishok--LeBaron and Moskowitz--Ooi--Pedersen}
\label{sec:08-lab}

\begin{callout}
\textbf{Lab~\ref{ch:08} at a glance.} \emph{Papers replicated:} Brock, Lakonishok and LeBaron (1992); Sullivan, Timmermann and White (1999); Moskowitz, Ooi and Pedersen (2012). \emph{Claim on trial\index{claims-on-trial method}:} simple trend rules beat buy-and-hold after correcting for the number of rules tried. \emph{Data:} the daily U.S. market index built from the Kenneth R. French Data Library\index{Kenneth French data library}, from July 1926, and optionally the daily Dow Jones Industrial Average from a free source such as Stooq; AQR's free \emph{Time Series Momentum: Factors, Monthly} data set; and, for a small multi-asset test, the 10-year Treasury yield (GS10), the H.10 exchange rates, three-month interbank rates and the WTI oil price from FRED\index{FRED}, which earlier labs already use. \emph{Code:} runs as a script or as a notebook, with a prediction prompt before each step.
\end{callout}

The lab has five parts. Each signal uses only prices up to the day, or month, before the position is taken.

\begin{itemize}
  \item \textbf{Part A, the 26 rules.} The lab applies Brock, Lakonishok and LeBaron's ten variable-length, ten fixed-length and six breakout rules to the daily U.S. market. For each rule it reports the number of buy and sell days, the mean daily return after buy and after sell signals, and the buy--sell difference with its $t$-statistic. It does this separately for 1926--1986, the overlap with the original sample, and for 1987 to the present, the period after publication.
  \item \textbf{Part B, the search correction.} The lab builds a universe of a few thousand moving-average and breakout rules by varying the windows, bands and holding periods. It computes each rule's excess return over buy-and-hold and applies White's Reality Check and Hansen's SPA test, using the stationary bootstrap with an average block length of ten days and 1,000 bootstrap samples. Again the two periods are reported separately.
  \item \textbf{Part C, time-series momentum across asset classes.} The lab reads AQR's monthly time-series momentum factors for equity indices, bonds, currencies and commodities, and for the diversified portfolio. It reports the average return, Sharpe ratio and alpha against the market, before and after the original sample ended in 2009.
  \item \textbf{Part D, the smile.} Using the AQR diversified factor and the U.S. market return, the lab regresses quarterly trend returns on the market's return and its square, as in Moskowitz, Ooi and Pedersen, and lists the strategy's return in the ten worst quarters for U.S. stocks.
  \item \textbf{Part E, asset by asset.} With free data the lab builds a small universe of about a dozen assets: U.S. stocks, a 10-year Treasury bond built from yields as in Lab~\ref{ch:13}, nine currencies against the dollar with their interest differentials as in Lab~\ref{ch:06}, and spot crude oil. It runs Huang and coauthors' tests on this universe: separate predictive regressions for each asset, the pooled regression with a bootstrap $p$-value, and a comparison of time-series momentum with the time-series history strategy.
\end{itemize}

\textbf{What readers will find.} Lab Figure~\ref{ch:08}.1 compares the buy--sell differences of the 26 rules before and after 1987. Lab Figure~\ref{ch:08}.2 shows the Reality Check and SPA $p$-values for the best rule in a rolling ten-year window. Lab Figure~\ref{ch:08}.3 plots quarterly trend returns against market returns, with the fitted smile. The test of the claim is whether the best rule survives the search correction after 1986, and whether the diversified strategy's return holds up after 2009 and asset by asset.

\emph{[Draft note: the figures and numbers are generated by running the lab on data at the time of publication. The book prints the results; the text above does not anticipate them.]}

\subsection{Robustness checks}
\label{sec:08-stress}

The full-sample results can hide weaknesses. The checks below test the claim under the conditions in which it is hardest to meet.

\textbf{1. After publication.} The cleanest test of a trading rule is the period after the paper that made it famous. For the stock-index rules that period begins in 1992, or in 1987 if one counts from the end of the sample. For time-series momentum it begins in 2010. Industry reports and the AQR data show that diversified trend strategies had a weak decade after 2009, with long flat periods, before a strong year in 2022 when inflation surged and bonds and stocks fell together. A strategy that works in one regime and stalls in another is not disproved by a flat decade, but the post-publication record is the evidence an investor should weigh most.

\textbf{2. Costs and a one-day lag.} The stock-index rules switch between long and short many times a year, and most of their profit comes from short holding periods. Bajgrowicz and Scaillet (2012) found that even low transaction costs offset the in-sample performance of the best rules on the Dow, and the lab measures how much a cost of 0.1\% per switch, or trading at the next day's close instead of today's, takes away. Time-series momentum with a twelve-month signal changes position far less often than daily rules, and it trades liquid futures, so its costs are much lower. Hurst, Ooi and Pedersen assumed costs six times higher before 1993 to reflect older markets and still found positive net returns.

\textbf{3. Asset by asset or pooled.} A pooled test can find a significant result when no single asset is significant, and Huang and coauthors argue that this is what happened. The answer depends on what the investor holds. An investor who trades a single market needs that market to be predictable. An investor who holds a diversified portfolio of fifty markets benefits from many weak signals combined, in the same way that Chapter~\ref{ch:01} showed a higher Sharpe ratio requires fewer years to detect. The lab reports both views.

\textbf{4. Is it trend or carry?} Here \emph{carry} means the return an investor earns simply by holding an asset that tends to pay off on average, as stocks and bonds have. The time-series history strategy of Huang and coauthors earns its return by holding assets that have had positive average returns, mainly stocks and bonds. If trend following's return over a long sample is mostly that, then its value lies elsewhere: in its tendency to turn short during long declines. Part D tests that tendency directly.

\textbf{5. Data snooping in the defense.} The one-, three- and twelve-month windows, the 40\% target and the 60-day volatility estimate are themselves choices. They were made by researchers who knew the history. The lab lets readers vary each choice and reports how sensitive the results are, in the spirit of the robustness grids in Chapter~\ref{ch:24}.

\subsection*{Your verdict}

Reach your own verdict on the claim from the numbers in your own run, using the scale of Section~\ref{sec:01-method}. Then compare it with the published evidence summarized in Section~\ref{sec:08-verdict}. If your verdict differs, explain which of your results led you there.

\section{Where the evidence stands}
\label{sec:08-verdict}

The published studies cited in this chapter point to a clear pattern: \textbf{classic trend rules on a single stock index do not survive a full accounting for the number of rules tried, while diversified trend following across many markets has a stronger record, though its average return is smaller and less certain than the original studies suggested.} The table summarizes where the evidence supports the claim and where it does not, as of November 2026, and names the source of each finding. The lab lets you verify these findings with free data.

\begin{center}
\small
\begin{tabularx}{\linewidth}{@{}>{\raggedright\arraybackslash}p{0.19\linewidth}>{\raggedright\arraybackslash}p{0.10\linewidth}>{\raggedright\arraybackslash}X>{\raggedright\arraybackslash}p{0.23\linewidth}@{}}
\toprule
Situation & Holds? & Why & Evidence \\
\midrule
Stock-index rules in the original sample & \textbf{Yes} & Returns after buy signals far exceeded returns after sell signals, and the best rules stayed significant even after the Reality Check. & Brock, Lakonishok and LeBaron (1992); Sullivan, Timmermann and White (1999) \\
\addlinespace[3pt]
Stock-index rules after 1986 & \textbf{No} & The best rule was no longer significant, and costs and a one-day lag removed much of the profit. & Sullivan, Timmermann and White (1999); Bajgrowicz and Scaillet (2012) \\
\addlinespace[3pt]
Time-series momentum across many markets & \textbf{Yes} & The strategy was profitable in all 58 contracts, with a Sharpe ratio above one for the diversified portfolio. & Moskowitz, Ooi and Pedersen (2012) \\
\addlinespace[3pt]
A century of trend following & \textbf{Yes} & A diversified trend strategy earned positive returns in every decade since 1880, after fees. & Hurst, Ooi and Pedersen (2017) \\
\addlinespace[3pt]
Asset by asset & \textbf{No} & Most single-asset slopes are insignificant, and the pooled test overstated the evidence. & Huang, Li, Wang and Zhou (2020) \\
\addlinespace[3pt]
After 2009 & \textbf{Unclear} & Diversified trend strategies had a long weak period, then strong gains in 2022. & AQR time-series momentum data \\
\addlinespace[3pt]
In long market declines & \textbf{Yes} & Trend following made money in most long equity drawdowns, which makes it a useful diversifier. & Moskowitz, Ooi and Pedersen (2012); Hurst, Ooi and Pedersen (2017) \\
\bottomrule
\end{tabularx}
\end{center}

The trend is not a reliable friend for someone who trades a single stock index with a classic moving-average rule. Those rules worked in the century in which they were found, and not after. Diversified trend following has a much stronger record, but part of its return may come from holding assets that simply rose over time, and its most dependable feature is its tendency to make money during long market declines.

\textbf{What would change the picture.} The evidence would support the weak form of the claim more strongly if diversified trend following kept earning a positive return after 2022 and if more single assets showed significant predictability on their own. It would weaken the claim further if the strategy's return turned out to come mainly from holding stocks and bonds on average rather than from trends.

\subsection{Open questions}
\label{sec:08-open}

\begin{itemize}
  \item \textbf{Has trend following become crowded?} Assets in managed-futures funds grew sharply after 2008. Did the extra money weaken trends, and is that why returns were lower after 2009?
  \item \textbf{Trend or carry?} How much of trend following's long-run return is predictability, and how much is the average return of the assets it tends to hold?
  \item \textbf{Which window?} Short windows react faster to turning points but trade more. Is there a principled way to choose, or to combine, look-back periods without overfitting?
  \item \textbf{Will the crisis behavior continue?} Trend following profited in long declines such as 2008 and 2022. Will it do so in the next crisis, which may be faster or shaped differently?
\end{itemize}

\section{What practitioners and investors can use}
\label{sec:08-investors}

The evidence in Section~\ref{sec:08-verdict} suggests several practical lessons.

\textbf{Use trend as a diversifier, not as a stand-alone bet.} A modest allocation to a diversified trend strategy has historically reduced the drawdowns of a stock--bond portfolio. It is not a strategy that wins every year, and it can go through years of flat or negative returns.

\textbf{Prefer breadth and patience.} A trend strategy that holds many markets in several asset classes and uses slow signals is more robust than one that trades a single market with fast signals. Fast rules generate more trades, higher costs and more room for overfitting.

\textbf{Check fees carefully.} Trend-following funds have often charged high fees. Simple rules on liquid futures are cheap to run, so the fee should be judged against what a transparent rule would deliver.

\textbf{Questions to ask a trend-following manager:}

\begin{enumerate}
  \item How many signal variants did you test before choosing the ones you use, and were the choices fixed before the backtest\index{backtest}?
  \item How did the strategy perform after 2009, and in 2022?
  \item What would the strategy have done in a one-day crash such as October 1987, and in a sharp rebound such as March 2009?
  \item How much of the long-run return comes from being long stocks and bonds on average?
  \item What are the turnover, trading costs and fees, and what is the return after all three?
\end{enumerate}

\begin{callout}
\textbf{What Failed: fast trend rules on a single stock index.} For decades, newsletters and trading systems sold moving-average rules for timing the stock market, often with backtests on the Dow going back to the early twentieth century. The rules that looked best were short-window rules that traded many times a year. Sullivan, Timmermann and White's test showed what happened next: from 1987 to 1996 the best rules were no longer significant, and on S\&P 500 futures, the market where the rules could actually be traded cheaply, no rule beat the benchmark. The failure was in treating the best rule in a long backtest as a rule with real predictive power. The claim that a fast moving-average rule on a single index beats buy-and-hold \textbf{fails}.
\end{callout}

\section*{Exercises}
\addcontentsline{toc}{section}{Exercises}

\subsection*{Conceptual (all readers)}

\begin{enumerate}
  \item An index is at 1,050. Its 200-day moving average is 1,030 and its 50-day high is 1,060. What signals do the 1--200 rule with a 1\% band and the 50-day breakout rule with a 1\% band give?
  \item In a year when every asset class rises, what does a time-series momentum strategy hold? What does the cross-sectional momentum strategy of Chapter~\ref{ch:05} hold?
  \item Explain why the payoff of a trend-following strategy resembles an option straddle. When does the analogy break down?
  \item Why does adding many poor rules to a universe make White's Reality Check less powerful, and how does Hansen's SPA test address the problem?
  \item Huang and coauthors' time-series history strategy needs no predictability. Explain how it can still earn a positive return.
  \item A newsletter claims that its moving-average signal ``would have avoided every major crash since 1950.'' Without any calculation, list the questions from this chapter you would ask before following its signals with your own savings.
\end{enumerate}

\subsection*{Lab extensions (technical readers)}

\begin{enumerate}
  \setcounter{enumi}{6}
  \item In Part A, charge 0.1\% for every switch between long and short and trade at the next day's close. How many of the 26 rules still beat buy-and-hold after 1986?
  \item In Part B, double the size of the rule universe by adding filter rules. How do the Reality Check and SPA $p$-values of the best rule change?
  \item In Part C, split the post-2009 period into 2010--2021 and 2022 to the present. How much of the post-publication return comes from 2022?
  \item In Part E, change the look-back from 12 months to 3 and to 24 months. Does any single asset become significant?
  \item Combine the AQR diversified trend factor with the U.S. market in an 80/20 portfolio. Compare its Sharpe ratio and maximum drawdown with the market alone, before and after 2009.
\end{enumerate}
